\SourceLeaf{354}%
varies from place to place on the surfaces. In \textsc{Mercator}'s projection, for example, the scale ratio is greater the farther the places in question are from the equator. In the stereographic projection it is greater the farther they are from the viewpoint.

There are infinitely many conformal mappings of any given surface onto another given surface. The general solution comprises them all by including an arbitrary function, which can be chosen at will or to suit the purpose at hand. If only part of one surface is to be mapped, it is generally best to choose a conformal mapping in which the variation in the scale ratio stays within the narrowest possible limits over the region to be mapped.

The problem of mapping the ellipsoid conformally onto the sphere is treated in its own right among the examples in the work cited. Two particular solutions are given alongside the general solution. One is especially suited to mapping the whole ellipsoid. The other is much better suited to the case in which only a moderate portion of the Earth's surface, regarded as ellipsoidal, is to be mapped conformally onto a sphere, as is always the case in specific applications to geodesy. An example given in that work shows how well this second mapping meets the requirement stated above. Within a zone spanning five degrees of latitude, the variation in the scale ratio is only~\(\frac1{530000}\). That work also sketches the main features of the method by which a conformal mapping can in general be used to calculate a triangulation network. The detailed execution and the application to this particular mapping were left for later treatment.

The present memoir is meant to fulfill this obligation, though in a somewhat different way from the one originally envisaged. It takes as its basis a third particular solution of the problem, different from the one just mentioned, which achieves the intended purpose even more fully. In these pages we must content ourselves with giving only a general idea of it.
